\chapter{Bindungsenergie und Weizsaecker-Massenformel}

\section{Ziel dieses Kapitels}

In diesem Kapitel wird erklaert, warum Atomkerne gebunden sind, warum die
Bindungsenergie pro Nukleon eine zentrale Groesse ist und wie die
Weizsaecker-Massenformel die wichtigsten Beitraege zur Kernmasse beschreibt.

\begin{notebox}
Dieses Kapitel behandelt die Stichwoerter:
\begin{itemize}
    \item Massendefekt,
    \item Bindungsenergie,
    \item Bindungsenergie pro Nukleon,
    \item Weizsaecker-Massenformel,
    \item Volumenterm,
    \item Oberflaechenterm,
    \item Coulombterm,
    \item Asymmetrieterm,
    \item Paarungsterm,
    \item Isobarenreihe,
    \item Massenparabel,
    \item Massental,
    \item Kernstabilitaet.
\end{itemize}
\end{notebox}

\section{Wiederholung: Massendefekt und Bindungsenergie}

Ein gebundener Atomkern ist leichter als die Summe seiner freien Nukleonen.
Diese Massendifferenz heisst Massendefekt.

\begin{definitionbox}
Der \emph{Massendefekt} $\Delta m$ eines Kerns ist
\[
    \Delta m = Zm_p + Nm_n - m_{\mathrm{Kern}}.
\]
Dabei ist:
\begin{itemize}
    \item $Z$ die Protonenzahl,
    \item $N$ die Neutronenzahl,
    \item $m_p$ die Masse eines freien Protons,
    \item $m_n$ die Masse eines freien Neutrons,
    \item $m_{\mathrm{Kern}}$ die Masse des gebundenen Kerns.
\end{itemize}
\end{definitionbox}

\begin{definitionbox}
Die \emph{Bindungsenergie} $B$ ist die Energie, die benoetigt wird, um einen Kern
vollstaendig in freie Protonen und Neutronen zu zerlegen.
Sie ist durch
\[
    B = \Delta m c^2
\]
gegeben.
\end{definitionbox}

\begin{formulabox}
\[
    B
    =
    \left(
        Zm_p + Nm_n - m_{\mathrm{Kern}}
    \right)c^2.
\]
\end{formulabox}

\begin{conceptbox}
Ein gebundener Kern hat eine niedrigere Energie als die getrennten Nukleonen.
Nach
\[
    E = mc^2
\]
bedeutet niedrigere Energie auch niedrigere Masse.

Die fehlende Masse ist also nicht verschwunden.
Sie entspricht der Bindungsenergie.
\end{conceptbox}

\section{Bindungsenergie pro Nukleon}

Die gesamte Bindungsenergie $B$ waechst grob mit der Groesse des Kerns.
Um Kerne verschiedener Groesse sinnvoll zu vergleichen, verwendet man die
Bindungsenergie pro Nukleon.

\begin{definitionbox}
Die \emph{Bindungsenergie pro Nukleon} ist
\[
    \frac{B}{A}.
\]
Sie gibt an, wie stark ein einzelnes Nukleon im Kern im Mittel gebunden ist.
\end{definitionbox}

\begin{conceptbox}
Die Vorlesungsabbildung zur Bindungsenergie pro Nukleon zeigt:
\begin{itemize}
    \item Bei sehr leichten Kernen steigt $B/A$ stark an.
    \item Bei mittleren Kernen erreicht $B/A$ ein Maximum.
    \item Bei sehr schweren Kernen faellt $B/A$ langsam wieder ab.
\end{itemize}

Damit erklaert diese Kurve qualitativ zwei wichtige Energiequellen:
\begin{itemize}
    \item Fusion leichter Kerne,
    \item Spaltung schwerer Kerne.
\end{itemize}
\end{conceptbox}

\begin{examtipbox}
Wenn eine Reaktion zu Produkten mit groesserer Bindungsenergie pro Nukleon fuehrt,
kann Energie frei werden.

Das ist der Grundgedanke hinter:
\[
    \text{leichte Kerne} \rightarrow \text{Fusion}
\]
und
\[
    \text{schwere Kerne} \rightarrow \text{Spaltung}.
\]
\end{examtipbox}

\section{Reaktionsenergie und Q-Wert}

Bei Kernreaktionen vergleicht man die Massen vor und nach der Reaktion.

\begin{definitionbox}
Der \emph{$Q$-Wert} einer Kernreaktion ist
\[
    Q
    =
    \left(
        m_{\mathrm{Anfang}} - m_{\mathrm{Ende}}
    \right)c^2.
\]
\end{definitionbox}

\begin{conceptbox}
Wenn
\[
    Q>0
\]
ist, wird Energie frei.
Die Reaktion ist exoenergetisch.

Wenn
\[
    Q<0
\]
ist, muss Energie zugefuehrt werden.
Die Reaktion ist endoenergetisch.
\end{conceptbox}

\begin{formulabox}
\[
    Q>0
    \quad\Longrightarrow\quad
    \text{Energie wird frei}.
\]
\[
    Q<0
    \quad\Longrightarrow\quad
    \text{Energie muss zugefuehrt werden}.
\]
\end{formulabox}

\section{Motivation fuer die Weizsaecker-Massenformel}

Die Weizsaecker-Massenformel beschreibt Kernmassen naeherungsweise.
Sie basiert auf der Idee, dass ein Kern in manchen Aspekten wie ein geladener
Fluessigkeitstropfen behandelt werden kann.

\begin{conceptbox}
Die Grundidee ist:
\begin{center}
\emph{Die Bindungsenergie eines Kerns setzt sich aus mehreren Beitraegen zusammen.}
\end{center}

Diese Beitraege beschreiben:
\begin{itemize}
    \item die starke Bindung im Kernvolumen,
    \item die fehlenden Nachbarn an der Oberflaeche,
    \item die Coulomb-Abstossung der Protonen,
    \item die energetische Strafe fuer $N\neq Z$,
    \item Paarungseffekte zwischen Nukleonen.
\end{itemize}
\end{conceptbox}

\begin{notebox}
In den Vorlesungsfolien wird die Bindungsenergie pro Nukleon in Beitraege zerlegt:
\begin{itemize}
    \item Kondensationsenergie,
    \item Oberflaechenenergie,
    \item Coulombenergie,
    \item Asymmetrie-Energie.
\end{itemize}
Diese Zerlegung entspricht qualitativ den Termen der Weizsaecker-Massenformel.
\end{notebox}

\section{Die Weizsaecker-Massenformel}

\begin{definitionbox}
Die \emph{Weizsaecker-Massenformel} beziehungsweise semi-empirische Massenformel
naeherert die Bindungsenergie eines Kerns mit Massenzahl $A$ und Protonenzahl $Z$.
\end{definitionbox}

\begin{formulabox}
Eine uebliche Schreibweise fuer die Bindungsenergie ist:
\[
    B(A,Z)
    =
    a_V A
    - a_S A^{2/3}
    - a_C \frac{Z(Z-1)}{A^{1/3}}
    - a_A \frac{(N-Z)^2}{A}
    + \delta(A,Z).
\]
Mit
\[
    N=A-Z
\]
kann man auch schreiben:
\[
    B(A,Z)
    =
    a_V A
    - a_S A^{2/3}
    - a_C \frac{Z(Z-1)}{A^{1/3}}
    - a_A \frac{(A-2Z)^2}{A}
    + \delta(A,Z).
\]
\end{formulabox}

\begin{notebox}
Man findet in Buechern auch Varianten mit
\[
    Z^2
\]
statt
\[
    Z(Z-1).
\]
Fuer grosse $Z$ ist der Unterschied klein.
Die physikalische Aussage bleibt dieselbe:
Der Coulombterm waechst mit der Zahl der Protonenpaare.
\end{notebox}

\begin{conceptbox}
Die Vorzeichen sind wichtig.

Positive Beitraege erhoehen die Bindungsenergie.
Negative Beitraege verringern die Bindungsenergie.

Die starke Wechselwirkung liefert den positiven Hauptbeitrag.
Oberflaeche, Coulombabstossung und Neutron-Proton-Asymmetrie reduzieren die Bindung.
\end{conceptbox}

\section{Volumenterm}

\begin{definitionbox}
Der \emph{Volumenterm} ist
\[
    B_V = a_V A.
\]
\end{definitionbox}

\begin{conceptbox}
Der Volumenterm beschreibt, dass jedes Nukleon im Inneren des Kerns durch die
starke Wechselwirkung mit seinen Nachbarn gebunden wird.

Da die starke Wechselwirkung kurzreichweitig ist, wechselwirkt ein Nukleon nur mit
einer begrenzten Zahl von Nachbarn.
Deshalb ist der Beitrag naeherungsweise proportional zur Zahl der Nukleonen:
\[
    B_V \propto A.
\]
\end{conceptbox}

\begin{examtipbox}
Der Volumenterm ist positiv:
\[
    +a_V A.
\]
Er ist der Hauptgrund, warum Kerne ueberhaupt gebunden sind.
\end{examtipbox}

\section{Oberflaechenterm}

\begin{definitionbox}
Der \emph{Oberflaechenterm} ist
\[
    B_S = -a_S A^{2/3}.
\]
\end{definitionbox}

\begin{conceptbox}
Nukleonen an der Kernoberflaeche haben weniger Nachbarn als Nukleonen im Inneren.
Sie sind deshalb weniger stark gebunden.

Da die Oberflaeche einer Kugel mit
\[
    R^2
\]
skaliert und
\[
    R \propto A^{1/3}
\]
gilt, skaliert die Oberflaeche wie
\[
    R^2 \propto A^{2/3}.
\]

Deshalb ist der Oberflaechenterm proportional zu
\[
    A^{2/3}.
\]
\end{conceptbox}

\begin{derivationbox}
Aus
\[
    R=r_0A^{1/3}
\]
folgt fuer die Oberflaeche:
\[
    S = 4\pi R^2
      = 4\pi r_0^2 A^{2/3}.
\]
Die Zahl der Nukleonen an der Oberflaeche ist daher naeherungsweise proportional zu
\[
    A^{2/3}.
\]
Da diese Nukleonen weniger stark gebunden sind, wird die Bindungsenergie um einen
Term
\[
    -a_SA^{2/3}
\]
verringert.
\end{derivationbox}

\begin{examtipbox}
Der Oberflaechenterm ist negativ:
\[
    -a_SA^{2/3}.
\]
Er ist besonders wichtig fuer leichte Kerne, weil dort der Anteil der
Oberflaechennukleonen gross ist.
\end{examtipbox}

\section{Coulombterm}

\begin{definitionbox}
Der \emph{Coulombterm} beschreibt die elektrische Abstossung zwischen den Protonen
im Kern.
\[
    B_C
    =
    -a_C \frac{Z(Z-1)}{A^{1/3}}.
\]
\end{definitionbox}

\begin{conceptbox}
Protonen tragen positive elektrische Ladung.
Sie stossen sich durch die elektromagnetische Wechselwirkung ab.

Diese Abstossung verringert die Bindungsenergie.
Deshalb hat der Coulombterm ein negatives Vorzeichen.
\end{conceptbox}

\begin{derivationbox}
Die Coulombenergie einer homogen geladenen Kugel skaliert grob wie
\[
    E_C \propto \frac{Z^2}{R}.
\]
Mit
\[
    R \propto A^{1/3}
\]
folgt
\[
    E_C \propto \frac{Z^2}{A^{1/3}}.
\]
Da Coulombabstossung die Bindung reduziert, erscheint dieser Beitrag in der
Bindungsenergie mit negativem Vorzeichen:
\[
    -a_C\frac{Z(Z-1)}{A^{1/3}}.
\]
\end{derivationbox}

\begin{conceptbox}
Der Coulombterm erklaert, warum sehr schwere Kerne weniger stark gebunden sind.
Je groesser $Z$ wird, desto staerker wird die elektrische Abstossung zwischen den
Protonen.
\end{conceptbox}

\section{Asymmetrieterm}

\begin{definitionbox}
Der \emph{Asymmetrieterm} bestraft ein Ungleichgewicht zwischen Neutronen und Protonen:
\[
    B_A
    =
    -a_A \frac{(N-Z)^2}{A}.
\]
\end{definitionbox}

\begin{conceptbox}
Der Asymmetrieterm ist quantenmechanisch motiviert.

Protonen und Neutronen sind Fermionen.
Wenn sehr viel mehr Neutronen als Protonen vorhanden sind, muessen Neutronen hoehere
Energiezustaende besetzen.
Das erhoeht die Gesamtenergie und verringert damit die Bindungsenergie.
\end{conceptbox}

\begin{formulabox}
Mit
\[
    A=N+Z
\]
und
\[
    N-Z = A-2Z
\]
kann man den Asymmetrieterm auch schreiben als
\[
    B_A
    =
    -a_A \frac{(A-2Z)^2}{A}.
\]
\end{formulabox}

\begin{derivationbox}
Die Vorlesungsherleitung des Asymmetrieterms benutzt eine Fermigasbetrachtung.

Fuer ein symmetrisches System mit
\[
    N=Z
\]
sind die Ferminiveaus von Neutronen und Protonen gleichmaessig gefuellt.

Betrachtet man nun einen asymmetrischen Kern mit
\[
    N\neq Z,
\]
dann ergibt sich fuer die kinetische Gesamtenergie ein Zusatzbeitrag gegenueber dem
symmetrischen Fall.

Mit
\[
    \Delta X = N-Z
\]
fuehrt die Entwicklung fuer kleine
\[
    \frac{\Delta X}{A}
\]
auf
\[
    \Delta E_{\mathrm{ges}}
    \propto
    A\left(\frac{\Delta X}{A}\right)^2.
\]
Damit:
\[
    \Delta E_{\mathrm{ges}}
    \propto
    \frac{(N-Z)^2}{A}.
\]
Da diese Zusatzenergie die Bindung verringert, erscheint der Asymmetrieterm in der
Bindungsenergie mit negativem Vorzeichen:
\[
    -a_A\frac{(N-Z)^2}{A}.
\]
\end{derivationbox}

\begin{examtipbox}
Der Asymmetrieterm erklaert, warum Kerne mit starkem Ungleichgewicht zwischen
Neutronen und Protonen weniger stabil sind.

Er bevorzugt naeherungsweise
\[
    N\approx Z.
\]
Bei schweren Kernen verschiebt der Coulombterm die Stabilitaet aber zu
\[
    N>Z.
\]
\end{examtipbox}

\section{Paarungsterm}

\begin{definitionbox}
Der \emph{Paarungsterm} $\delta(A,Z)$ beschreibt, dass Kerne mit gepaarten Protonen
und gepaarten Neutronen besonders stark gebunden sind.
\end{definitionbox}

\begin{conceptbox}
Nukleonen bilden bevorzugt Paare mit entgegengesetztem Spin.

Daraus folgt:
\begin{itemize}
    \item gerade-gerade Kerne sind besonders stabil,
    \item ungerade-ungerade Kerne sind weniger stabil,
    \item bei ungeradem $A$ ist der Paarungseffekt oft naeherungsweise null.
\end{itemize}
\end{conceptbox}

\begin{formulabox}
Eine uebliche Schreibweise ist:
\[
    \delta(A,Z)
    =
    \begin{cases}
        +a_P A^{-1/2}, & \text{gerade } Z \text{ und gerade } N,\\
        0, & \text{ungerades } A,\\
        -a_P A^{-1/2}, & \text{ungerade } Z \text{ und ungerade } N.
    \end{cases}
\]
\end{formulabox}

\begin{examtipbox}
Der Paarungsterm ist der Grund, warum bei geradem $A$ zwei Massenparabeln auftreten
koennen:
\begin{itemize}
    \item eine fuer gerade-gerade Kerne,
    \item eine fuer ungerade-ungerade Kerne.
\end{itemize}
\end{examtipbox}

\section{Gesamte Formel fuer die Kernmasse}

Die Weizsaecker-Formel wird oft als Formel fuer die Bindungsenergie angegeben.
Die Kernmasse ergibt sich dann aus den freien Nukleonenmassen minus Bindungsenergie.

\begin{formulabox}
\[
    m_{\mathrm{Kern}}(A,Z)c^2
    =
    Zm_pc^2
    +
    Nm_nc^2
    -
    B(A,Z).
\]
\end{formulabox}

\begin{conceptbox}
Je groesser die Bindungsenergie $B(A,Z)$ ist, desto kleiner ist die Kernmasse.

Deshalb entspricht maximale Bindungsenergie bei festem $A$ einem Massenminimum.
Dieses Massenminimum ist entscheidend fuer die Stabilitaet innerhalb einer
Isobarenreihe.
\end{conceptbox}

\section{Isobarenreihe und Massenparabel}

Bei einer Isobarenreihe ist die Massenzahl
\[
    A
\]
fest.
Dann gilt:
\[
    N=A-Z.
\]
Die Masse wird also als Funktion von $Z$ betrachtet.

\begin{definitionbox}
Eine \emph{Isobarenreihe} ist eine Reihe von Kernen mit gleicher Massenzahl $A$,
aber unterschiedlicher Protonenzahl $Z$.
\end{definitionbox}

\begin{conceptbox}
Setzt man
\[
    N=A-Z
\]
in die Weizsaecker-Massenformel ein und betrachtet festes $A$, dann haengt die Masse
hauptsaechlich quadratisch von $Z$ ab.

Dadurch entsteht eine \emph{Massenparabel}.
Der stabilste Kern einer Isobarenreihe liegt am Massenminimum.
\end{conceptbox}

\begin{examplebox}
In den Vorlesungsfolien wird als Beispiel die Isobarenreihe mit
\[
    A=81
\]
gezeigt:
\[
    {}^{81}_{33}\mathrm{As},
    \quad
    {}^{81}_{34}\mathrm{Se},
    \quad
    {}^{81}_{35}\mathrm{Br},
    \quad
    {}^{81}_{36}\mathrm{Kr},
    \quad
    {}^{81}_{37}\mathrm{Rb}.
\]
Die Bindungsenergien pro Nukleon sind nahe beieinander, aber
\[
    {}^{81}_{35}\mathrm{Br}
\]
hat in der Tabelle den groessten Wert von $B/A$ und gleichzeitig die kleinste Masse.
\end{examplebox}

\begin{conceptbox}
Bei ungeradem $A$ gibt es typischerweise eine einzige Massenparabel.
Der Kern am Minimum ist stabil gegenueber Beta-Zerfall.

Kerne links und rechts vom Minimum koennen durch
\[
    \beta^-
\]
oder
\[
    \beta^+
\]
beziehungsweise Elektroneneinfang in Richtung des Minimums zerfallen.
\end{conceptbox}

\begin{examtipbox}
Merke:
\[
    \text{fester } A
    \quad\Rightarrow\quad
    \text{Isobarenreihe}.
\]
\[
    \text{Masse als Funktion von } Z
    \quad\Rightarrow\quad
    \text{Massenparabel}.
\]
\[
    \text{Minimum der Masse}
    \quad\Rightarrow\quad
    \text{maximale Stabilitaet}.
\]
\end{examtipbox}

\section{Beta-Zerfall und Richtung zum Massenminimum}

Die Massenparabel erklaert qualitativ, warum instabile Isobare durch Beta-Zerfall
in Richtung des Massenminimums gehen.

\begin{conceptbox}
Bei
\[
    \beta^-
\text{-Zerfall}
\]
wird ein Neutron in ein Proton umgewandelt:
\[
    n \rightarrow p + e^- + \overline{\nu}_e.
\]
Dadurch steigt $Z$ um eins, waehrend $A$ gleich bleibt.

Bei
\[
    \beta^+
\text{-Zerfall}
\]
wird ein Proton in ein Neutron umgewandelt:
\[
    p \rightarrow n + e^+ + \nu_e.
\]
Dadurch sinkt $Z$ um eins, waehrend $A$ gleich bleibt.
\end{conceptbox}

\begin{formulabox}
Fuer Beta-Zerfaelle gilt:
\[
    A=\text{konstant}.
\]
Beim $\beta^-$-Zerfall:
\[
    Z \rightarrow Z+1.
\]
Beim $\beta^+$-Zerfall:
\[
    Z \rightarrow Z-1.
\]
\end{formulabox}

\begin{examtipbox}
In einer Massenparabel zerfallen Kerne energetisch nach unten.
Der stabile Kern liegt am Minimum.

Links vom Minimum ist meist $\beta^-$-Zerfall relevant.
Rechts vom Minimum sind $\beta^+$-Zerfall oder Elektroneneinfang relevant.
\end{examtipbox}

\section{Massental}

Betrachtet man nicht nur eine feste Massenzahl $A$, sondern viele Kerne im
$(N,Z)$-Diagramm, dann ergibt sich ein Stabilitaetsbereich.

\begin{definitionbox}
Das \emph{Massental} ist die gedachte Landschaft der Kernmassen beziehungsweise
Bindungsenergien als Funktion von Neutronenzahl $N$ und Protonenzahl $Z$.

Stabile Kerne liegen entlang eines Tals beziehungsweise Stabilitaetsbandes.
\end{definitionbox}

\begin{conceptbox}
In den Vorlesungsabbildungen zur Bindungsenergie in Abhaengigkeit von $Z$ und $N$
ragt ein Stabilitaetsgebirge beziehungsweise ein Stabilitaetsband hervor.

Die stabilen Kerne liegen dort, wo die Bindungsenergie besonders gross beziehungsweise
die Masse besonders klein ist.
\end{conceptbox}

\begin{notebox}
In derselben Darstellung erscheinen auch magische Zahlen und ein Hinweis auf eine
moegliche Stabilitaetsinsel bei superschweren Kernen.

Die genaue Behandlung der magischen Zahlen gehoert zum Kapitel ueber Kernmodelle und
Schalenmodell.
\end{notebox}

\section{Warum schwere stabile Kerne neutronenreich sind}

Aus Kapitel 1 wissen wir:
Fuer schwere stabile Kerne gilt meist
\[
    N>Z.
\]
Die Weizsaecker-Massenformel erklaert diesen Trend.

\begin{conceptbox}
Der Asymmetrieterm bevorzugt
\[
    N\approx Z.
\]
Der Coulombterm bestraft aber viele Protonen, weil sich Protonen elektrisch abstossen.

Bei schweren Kernen wird die Coulombabstossung stark.
Um bei grossem $A$ stabil zu bleiben, ist es energetisch guenstiger, mehr Neutronen
als Protonen zu besitzen.
\end{conceptbox}

\begin{examtipbox}
Bei leichten Kernen gilt fuer stabile Kerne oft naeherungsweise
\[
    N\approx Z.
\]
Bei schweren Kernen gilt fuer stabile Kerne meist
\[
    N>Z.
\]
Der Grund ist die zunehmende Coulombabstossung der Protonen.
\end{examtipbox}

\section{Grenzen der Weizsaecker-Massenformel}

\begin{conceptbox}
Die Weizsaecker-Massenformel beschreibt viele globale Trends gut:
\begin{itemize}
    \item Bindungsenergie pro Nukleon,
    \item Massenparabeln,
    \item Stabilitaetslinie,
    \item ungefaehre Energiegewinne bei Spaltung und Fusion.
\end{itemize}
\end{conceptbox}

\begin{notebox}
Die Formel ist aber nur semi-empirisch.
Sie beschreibt nicht alle Feinstrukturen.

Insbesondere benoetigt man fuer spezielle Stabilitaeten:
\begin{itemize}
    \item Schalenmodell,
    \item magische Zahlen,
    \item detaillierte Kernstruktur,
    \item Deformationen.
\end{itemize}
\end{notebox}

\section{Mini-Zusammenfassung}

\begin{notebox}
Die wichtigsten Punkte dieses Kapitels sind:
\begin{itemize}
    \item Der Massendefekt ist
    \[
        \Delta m = Zm_p + Nm_n - m_{\mathrm{Kern}}.
    \]
    \item Die Bindungsenergie ist
    \[
        B=\Delta mc^2.
    \]
    \item Die Bindungsenergie pro Nukleon $B/A$ ist wichtig zum Vergleich verschiedener Kerne.
    \item Die Weizsaecker-Massenformel zerlegt die Bindungsenergie in mehrere Beitraege:
    \[
        B
        =
        a_V A
        - a_S A^{2/3}
        - a_C \frac{Z(Z-1)}{A^{1/3}}
        - a_A \frac{(N-Z)^2}{A}
        + \delta(A,Z).
    \]
    \item Der Volumenterm bindet.
    \item Der Oberflaechenterm reduziert die Bindung.
    \item Der Coulombterm reduziert die Bindung durch Proton-Proton-Abstossung.
    \item Der Asymmetrieterm bestraft $N\neq Z$.
    \item Der Paarungsterm erklaert besondere Stabilitaet von gerade-gerade Kernen.
    \item Bei festem $A$ entstehen Massenparabeln.
    \item Stabile Isobare liegen am Massenminimum.
    \item Das Massental beschreibt den Stabilitaetsbereich im $(N,Z)$-Diagramm.
\end{itemize}
\end{notebox}

\section{Pruefungscheck}

\begin{examtipbox}
Nach diesem Kapitel solltest du folgende Fragen beantworten koennen:
\begin{enumerate}
    \item Was ist der Massendefekt?
    \item Wie haengt Bindungsenergie mit Massendefekt zusammen?
    \item Warum betrachtet man $B/A$?
    \item Warum kann Fusion leichter Kerne Energie freisetzen?
    \item Warum kann Spaltung schwerer Kerne Energie freisetzen?
    \item Wie lautet die Weizsaecker-Massenformel?
    \item Was beschreibt der Volumenterm?
    \item Warum ist der Oberflaechenterm negativ?
    \item Warum ist der Coulombterm negativ?
    \item Warum skaliert der Coulombterm etwa wie $Z^2/A^{1/3}$?
    \item Was beschreibt der Asymmetrieterm?
    \item Warum hat der Asymmetrieterm die Form $(N-Z)^2/A$?
    \item Was beschreibt der Paarungsterm?
    \item Was ist eine Isobarenreihe?
    \item Warum entsteht bei festem $A$ eine Massenparabel?
    \item Warum liegt der stabile Isobar am Massenminimum?
    \item Warum sind schwere stabile Kerne neutronenreich?
    \item Was ist das Massental?
    \item Warum braucht man trotz Weizsaecker-Massenformel noch das Schalenmodell?
\end{enumerate}
\end{examtipbox}